\documentclass[../main.tex]{subfiles}
\begin{document}
% Việt hoá nhãn thuật toán (algorithm2e) cho nhất quán: "Thuật toán" thay vì "Algorithm"
\renewcommand{\algorithmcfname}{Thuật toán}

Đây là chương trọng tâm của đồ án. Chương trình bày thiết kế tổng thể và chi tiết của hệ
thống (kiến trúc, cơ sở dữ liệu, luồng xác thực, các thuật toán cốt lõi, quy trình điểm
danh và chống gian lận), tiếp đó là quá trình triển khai thành sản phẩm thực tế, và cuối
cùng là kết quả đánh giá trên các tiêu chí chức năng, bảo mật và hiệu năng.

% ═══════════════════════════════════════════════════════════════
% PHẦN A — PHÂN TÍCH VÀ THIẾT KẾ
% ═══════════════════════════════════════════════════════════════

\section{Thiết kế kiến trúc}
\label{sec:kien-truc}

Hệ thống được tổ chức theo mô hình ba tầng (three-tier), gồm tầng giao diện phía người
dùng, tầng xử lý nghiệp vụ và tầng lưu trữ dữ liệu, như minh hoạ trong
Hình~\ref{fig:kientruc}.

\begin{itemize}[leftmargin=1.5em]
  \item \textbf{Tầng giao diện (Client).} Bao gồm \textit{Zalo Mini App} phục vụ
  sinh viên cùng giảng viên và \textit{trang quản trị web} phục vụ phòng đào tạo.
  Tầng này chịu trách nhiệm hiển thị, thu nhận thao tác người dùng, chụp ảnh khuôn mặt và
  quét mã QR.

  \item \textbf{Tầng xử lý (Backend).} Gồm các hàm Cloud Functions xử lý các nghiệp vụ
  nhạy cảm: ký và xác minh mã QR, tính điểm tin cậy, phát hiện gian lận, xác thực và phân
  quyền. Việc đặt các logic quan trọng ở tầng máy chủ giúp tránh để lộ khoá bí mật hay bị
  can thiệp từ phía thiết bị người dùng.

  \item \textbf{Tầng lưu trữ (Storage).} Gồm cơ sở dữ liệu Firestore lưu dữ liệu nghiệp
  vụ và đồng bộ theo thời gian thực, cùng Firebase Storage lưu trữ dữ liệu ảnh.
\end{itemize}

\begin{figure}[H]
  \centering
  \includegraphics[width=0.72\textwidth]{kientruc}
  \caption{Kiến trúc tổng thể ba tầng của hệ thống}
  \label{fig:kientruc}
\end{figure}

\section{Thiết kế chi tiết}

Phần này trình bày thiết kế chi tiết của các thành phần cốt lõi trong hệ thống.

\subsection{Thiết kế cơ sở dữ liệu Firestore}
\label{sec:db-design}

Khác với cơ sở dữ liệu quan hệ, Firestore lưu dữ liệu theo mô hình tài liệu (document):
dữ liệu được tổ chức thành các \textit{collection}, mỗi collection chứa nhiều
\textit{document}. Hệ thống sử dụng các collection chính được liệt kê trong
Bảng~\ref{tab:collections}.

\begin{table}[H]
\centering
\small
\caption{Các collection chính của cơ sở dữ liệu}
\label{tab:collections}
\begin{tabular}{|l|p{9.5cm}|}
\hline
\textbf{Collection} & \textbf{Mô tả} \\
\hline
\texttt{users} & Hồ sơ người dùng (sinh viên, giảng viên, quản trị viên) \\ \hline
\texttt{classes} & Thông tin lớp học phần, lịch học, danh sách sinh viên \\ \hline
\texttt{sessions} & Phiên điểm danh, kèm khoá bí mật HMAC và cấu hình \\ \hline
\texttt{attendance} & Bản ghi điểm danh của từng sinh viên trong một phiên \\ \hline
\texttt{face\_registrations} & Dữ liệu khuôn mặt đã đăng ký của sinh viên \\ \hline
\texttt{pairing\_tokens} & Mã ghép nối cho tính năng Cổng máy chiếu \\ \hline
\texttt{fraud\_reports} & Báo cáo phát hiện gian lận theo phiên \\ \hline
\texttt{absence\_requests} & Đơn xin nghỉ của sinh viên \\ \hline
\end{tabular}
\end{table}

Thiết kế chi tiết của hai collection quan trọng nhất --- \texttt{sessions} và
\texttt{attendance} --- được trình bày trong Bảng~\ref{tab:session} và
Bảng~\ref{tab:attendance}. Sơ đồ thực thể --- liên kết tổng thể được trình bày trong
Hình~\ref{fig:erd}.

\begin{table}[H]
\centering
\small
\caption{Cấu trúc collection \texttt{sessions}}
\label{tab:session}
\begin{tabular}{|l|l|p{7.5cm}|}
\hline
\textbf{Trường} & \textbf{Kiểu} & \textbf{Ý nghĩa} \\
\hline
\texttt{classId} & string & Mã lớp học phần \\ \hline
\texttt{status} & string & Trạng thái: \texttt{active} / \texttt{ended} \\ \hline
\texttt{hmacSecret} & string & Khoá bí mật để ký mã QR \\ \hline
\texttt{qrRefreshInterval} & number & Chu kỳ làm mới mã QR (giây) \\ \hline
\texttt{faceRequired} & boolean & Có yêu cầu xác minh khuôn mặt không \\ \hline
\texttt{peerRequired} & boolean & Có yêu cầu xác minh ngang hàng không \\ \hline
\texttt{geoFenceRadius} & number & Bán kính hàng rào địa lý (mét) \\ \hline
\texttt{startedAt} & number & Thời điểm bắt đầu phiên \\ \hline
\end{tabular}
\end{table}

\begin{table}[H]
\centering
\small
\caption{Cấu trúc collection \texttt{attendance}}
\label{tab:attendance}
\begin{tabular}{|l|l|p{7.2cm}|}
\hline
\textbf{Trường} & \textbf{Kiểu} & \textbf{Ý nghĩa} \\
\hline
\texttt{sessionId} & string & Mã phiên điểm danh \\ \hline
\texttt{studentId} & string & Mã sinh viên \\ \hline
\texttt{peerVerifications} & array & Danh sách xác minh ngang hàng \\ \hline
\texttt{peerCount} & number & Số lượt xác minh ngang hàng \\ \hline
\texttt{trustPolicy} & string & Điểm tin cậy: \texttt{present}/\texttt{review}/\texttt{absent} \\ \hline
\texttt{faceVerification} & object & Kết quả xác minh khuôn mặt \\ \hline
\texttt{location} & object & Toạ độ GPS lúc điểm danh \\ \hline
\end{tabular}
\end{table}

Ngoài hai collection trên, cấu trúc của một số collection quan trọng khác được trình bày
trong Bảng~\ref{tab:users} và Bảng~\ref{tab:classes}.

\begin{table}[H]
\centering
\small
\caption{Cấu trúc collection \texttt{users}}
\label{tab:users}
\begin{tabular}{|l|l|p{7.3cm}|}
\hline
\textbf{Trường} & \textbf{Kiểu} & \textbf{Ý nghĩa} \\
\hline
\texttt{id} & string & Định danh người dùng (theo Zalo) \\ \hline
\texttt{role} & string & Vai trò: \texttt{student}/\texttt{teacher}/\texttt{admin} \\ \hline
\texttt{name} & string & Họ và tên \\ \hline
\texttt{mssv} & string & Mã số sinh viên (nếu là sinh viên) \\ \hline
\texttt{faceRegistered} & boolean & Đã đăng ký khuôn mặt hay chưa \\ \hline
\texttt{hustVerified} & boolean & Đã xác minh email do trường cấp \\ \hline
\end{tabular}
\end{table}

\begin{table}[H]
\centering
\small
\caption{Cấu trúc collection \texttt{classes}}
\label{tab:classes}
\begin{tabular}{|l|l|p{7.3cm}|}
\hline
\textbf{Trường} & \textbf{Kiểu} & \textbf{Ý nghĩa} \\
\hline
\texttt{id} & string & Định danh lớp học phần \\ \hline
\texttt{code} & string & Mã lớp để sinh viên tham gia \\ \hline
\texttt{teacherId} & string & Định danh giảng viên phụ trách \\ \hline
\texttt{studentIds} & array & Danh sách sinh viên trong lớp \\ \hline
\texttt{schedule} & object & Lịch học (thứ, giờ bắt đầu, giờ kết thúc) \\ \hline
\end{tabular}
\end{table}

\begin{figure}[H]
  \centering
  \includegraphics[width=0.85\textwidth]{erd}
  \caption{Sơ đồ thực thể --- liên kết của cơ sở dữ liệu}
  \label{fig:erd}
\end{figure}

Vòng đời của phiên điểm danh và trạng thái của bản ghi điểm danh được mô tả bằng lược đồ
trạng thái trong Hình~\ref{fig:state}.

\begin{figure}[H]
  \centering
  \includegraphics[width=0.85\textwidth]{state_lifecycle}
  \caption{Lược đồ trạng thái của phiên và bản ghi điểm danh}
  \label{fig:state}
\end{figure}

\subsection{Thiết kế luồng xác thực và phân quyền}
\label{sec:auth-design}

Luồng xác thực dựa trên việc đăng nhập bằng tài khoản Zalo để lấy danh tính cơ bản, sau đó
sinh Custom Token của Firebase để thống nhất danh tính người dùng trên toàn hệ thống.

Việc phân quyền được thực thi ở tầng cơ sở dữ liệu bằng quy tắc bảo mật (Security Rules)
của Firestore. Một nguyên tắc quan trọng là \textbf{ngăn người dùng tự nâng quyền}: tài
khoản tạo từ phía thiết bị chỉ được phép mang vai trò \texttt{student}, còn vai trò
\texttt{teacher} và \texttt{admin} chỉ được gán bởi Cloud Function. Đoạn quy tắc minh hoạ
cho collection \texttt{users} như sau:

\begin{lstlisting}[caption={Quy tắc bảo mật cho collection users},captionpos=b]
match /users/{userId} {
  allow read:   if true;
  allow create: if newDocSafeRole();   // role chi duoc la "student"
  allow update: if notChangingRole();  // khong duoc tu doi role
  allow delete: if false;
}
\end{lstlisting}

Tương tự, collection \texttt{attendance} bị khoá ghi từ phía thiết bị --- mọi thao tác ghi
chỉ được thực hiện qua Cloud Function, như Danh sách mã nguồn~\ref{lst:attrule}.

\begin{lstlisting}[caption={Quy tắc bảo mật cho collection attendance},captionpos=b,label=lst:attrule]
match /attendance/{id} {
  allow read:   if true;
  allow create: if false;   // chi tao qua Cloud Function
  allow update: if false;
  allow delete: if false;
}
\end{lstlisting}

Ma trận phân quyền tổng quát theo vai trò được tóm tắt trong Bảng~\ref{tab:phanquyen}.

\begin{table}[H]
\centering
\small
\caption{Ma trận phân quyền theo vai trò}
\label{tab:phanquyen}
\begin{tabular}{|p{5cm}|c|c|c|}
\hline
\textbf{Thao tác} & \textbf{Sinh viên} & \textbf{Giảng viên} & \textbf{Quản trị} \\
\hline
Điểm danh & Có & --- & --- \\ \hline
Tạo/quản lý phiên điểm danh & --- & Có & --- \\ \hline
Duyệt, điểm danh thủ công & --- & Có & --- \\ \hline
Quản lý người dùng, lớp & --- & --- & Có \\ \hline
Duyệt đơn xin nghỉ & --- & --- & Có \\ \hline
\end{tabular}
\end{table}

\subsection{Thuật toán mã QR HMAC xoay 30 giây}
\label{sec:hmac-rotating}

Để chống việc chụp màn hình và chia sẻ lại mã QR, hệ thống sinh mã QR động được ký bằng
HMAC-SHA256 (xem công thức~(\ref{eq:hmac})) và làm mới sau mỗi 30 giây. Nội dung được ký
là chuỗi ghép từ năm trường: loại mã, mã phiên, mã người dùng, mốc thời gian và một giá
trị ngẫu nhiên (nonce). Thuật toán sinh mã được trình bày trong Thuật toán~\ref{algo:qr}.

\begin{algorithm}[H]
\caption{Sinh mã QR động của giảng viên}
\label{algo:qr}
\KwIn{Loại mã $type$, mã phiên $sid$, mã người dùng $uid$, khoá bí mật $K$}
\KwOut{Chuỗi mã QR đã ký}
$ts \leftarrow$ mốc thời gian hiện tại\;
$nonce \leftarrow$ chuỗi ngẫu nhiên 16 byte\;
$msg \leftarrow type : sid : uid : ts : nonce$\;
$sig \leftarrow \text{HMAC-SHA256}(msg, K)$\;
\Return $\text{JSON}(type, sid, uid, ts, nonce, sig)$\;
\end{algorithm}

Đoạn mã rút gọn cho việc sinh mã QR đã ký bằng HMAC được trình bày trong
Danh sách mã nguồn~\ref{lst:qrsign}.

\begin{lstlisting}[caption={Sinh mã QR đã ký HMAC-SHA256},captionpos=b,label=lst:qrsign]
function buildQR(type, sessionId, userId, secret) {
  const ts = Date.now();
  const nonce = randomHex(16);
  const msg = [type, sessionId, userId, ts, nonce].join(":");
  const sig = HmacSHA256(msg, secret).toString();
  return JSON.stringify({ type, sessionId, userId, ts, nonce, sig });
}
\end{lstlisting}

Phía sinh viên (hoặc máy chủ) xác minh mã bằng cách tính lại chữ ký từ cùng nội dung và
khoá bí mật rồi so sánh với chữ ký kèm theo; mã chỉ hợp lệ nếu chữ ký trùng khớp
\textbf{và} mốc thời gian còn trong khoảng dung sai 90 giây. Khoảng dung sai 90 giây
(gấp ba chu kỳ làm mới) nhằm dung hoà sai lệch đồng hồ giữa các thiết bị mà vẫn giới hạn
thời gian sống của mã. Nhờ đó, ngay cả khi mã QR bị chụp lại, nó cũng nhanh chóng hết
hạn và không thể dùng cho buổi học khác.

Toàn bộ quá trình khởi tạo khoá bí mật, sinh mã QR động và xác minh phía máy chủ được minh
hoạ trong Hình~\ref{fig:seq-qr}.

\begin{figure}[H]
  \centering
  \includegraphics[width=0.85\textwidth]{seq_qr}
  \caption{Lược đồ tuần tự sinh và xác minh mã QR HMAC}
  \label{fig:seq-qr}
\end{figure}

\subsection{Thiết kế giao diện và bộ nhận diện}
\label{sec:ui-design}

Giao diện ứng dụng tuân theo bộ nhận diện của Đại học Bách Khoa Hà Nội, với màu chủ đạo
là đỏ HUST (mã màu \texttt{\#be1d2c}). Hệ thống xây dựng một bộ thành phần giao diện dùng
chung nhằm bảo đảm tính nhất quán, tiêu biểu gồm: thẻ thống kê, vòng tròn điểm tin cậy
(score ring), thanh tiến trình, hộp thoại dạng trượt từ dưới lên, và trình quét mã QR
nhúng trực tiếp. Toàn bộ giao diện được tối ưu cho màn hình điện thoại với thao tác tối
giản, phù hợp với kỳ vọng về tốc độ của sinh viên.

\subsection{Luồng điểm danh bốn bước cho sinh viên}
\label{sec:attendance-flow}

Quy trình điểm danh của sinh viên gồm bốn bước tuần tự:

\begin{enumerate}[leftmargin=1.5em]
  \item \textbf{Quét mã QR của giảng viên.} Sinh viên dùng máy ảnh quét mã QR động đang
  hiển thị; hệ thống xác minh chữ ký HMAC và mốc thời gian.

  \item \textbf{Xác minh khuôn mặt.} Ứng dụng tự động chụp và so khớp khuôn mặt với ảnh
  đã đăng ký; kết quả đạt nếu độ tin cậy không nhỏ hơn ngưỡng 0{,}7.

  \item \textbf{Xác minh ngang hàng.} Sinh viên hiển thị mã QR cá nhân và quét mã của các
  bạn xung quanh; yêu cầu tối thiểu ba lượt xác minh.

  \item \textbf{Hoàn tất.} Hệ thống tổng hợp kết quả và tính điểm tin cậy.
\end{enumerate}

Quy trình điểm danh tổng quát cùng các điểm rẽ nhánh quyết định được minh hoạ dưới dạng lược
đồ hoạt động trong Hình~\ref{fig:activity}.

\begin{figure}[H]
  \centering
  \includegraphics[width=0.55\textwidth]{activity_diemdanh}
  \caption{Lược đồ hoạt động của quy trình điểm danh}
  \label{fig:activity}
\end{figure}

Điểm tin cậy (trust policy) được tính theo quy tắc: đạt \texttt{present} nếu vượt qua cả
điều kiện khuôn mặt và điều kiện ngang hàng; \texttt{review} (cần xem xét) nếu chỉ đạt
một trong hai; và \texttt{absent} nếu không đạt cả hai. Quy tắc này cho phép giảng viên
tập trung kiểm tra thủ công các trường hợp ở trạng thái \texttt{review}. Quy tắc tính điểm
tin cậy được trình bày hình thức trong Thuật toán~\ref{algo:trust}.

\begin{algorithm}[H]
\caption{Tính điểm tin cậy của một lượt điểm danh}
\label{algo:trust}
\KwIn{Cờ yêu cầu khuôn mặt $faceReq$, cờ yêu cầu ngang hàng $peerReq$, kết quả khuôn mặt $face$, số lượt ngang hàng $peerCount$}
\KwOut{Điểm tin cậy thuộc \{present, review, absent\}}
$facePass \leftarrow (\lnot faceReq) \lor face.matched$\;
$peerPass \leftarrow (\lnot peerReq) \lor (peerCount \geq 3)$\;
\eIf{$facePass$ \textbf{và} $peerPass$}{
  \Return present\;
}{
  \eIf{$facePass$ \textbf{hoặc} $peerPass$}{
    \Return review\;
  }{
    \Return absent\;
  }
}
\end{algorithm}

Toàn bộ trình tự tương tác của quy trình điểm danh bốn bước được minh hoạ trong
Hình~\ref{fig:seq-attendance}.

\begin{figure}[H]
  \centering
  \includegraphics[width=0.78\textwidth]{seq_attendance}
  \caption{Lược đồ tuần tự quy trình điểm danh bốn bước}
  \label{fig:seq-attendance}
\end{figure}

Quy trình đăng ký khuôn mặt --- điều kiện tiên quyết cho bước xác minh khi điểm danh --- được
minh hoạ trong Hình~\ref{fig:seq-face}.

\begin{figure}[H]
  \centering
  \includegraphics[width=0.95\textwidth]{seq_face}
  \caption{Lược đồ tuần tự đăng ký khuôn mặt}
  \label{fig:seq-face}
\end{figure}

\subsection{Hệ thống chống gian lận nhiều lớp}
\label{sec:anti-fraud}

Hệ thống áp dụng nguyên lý phòng thủ theo chiều sâu với nhiều lớp xác minh bổ trợ: mã QR
động (chống chia sẻ mã), xác minh khuôn mặt (chống điểm danh hộ), xác minh ngang hàng
(chống điểm danh từ xa) và kiểm tra vị trí GPS (chống điểm danh ngoài lớp). Bổ sung cho
các lớp này, một hàm phân tích chạy phía máy chủ rà soát các mẫu hành vi đáng ngờ và lập
báo cáo gian lận. Các loại hành vi đáng ngờ được hệ thống phát hiện được tổng hợp trong
Bảng~\ref{tab:fraud}.

\begin{table}[H]
\centering
\small
\caption{Phân loại các hành vi đáng ngờ}
\label{tab:fraud}
\begin{tabular}{|l|p{9cm}|}
\hline
\textbf{Loại} & \textbf{Mô tả} \\
\hline
\texttt{always\_same\_peers} & Luôn xác minh ngang hàng với cùng một nhóm bạn \\ \hline
\texttt{rapid\_verification} & Tốc độ xác minh nhanh bất thường \\ \hline
\texttt{low\_peer\_count} & Số lượt xác minh ngang hàng quá thấp \\ \hline
\texttt{face\_mismatch} & Khuôn mặt không khớp \\ \hline
\texttt{ai\_detected} & Mẫu bất thường do mô hình phân tích phát hiện \\ \hline
\end{tabular}
\end{table}

Thuật toán phân tích phát hiện gian lận theo phiên được trình bày trong
Thuật toán~\ref{algo:fraud}.

\begin{algorithm}[H]
\caption{Phân tích phát hiện gian lận theo phiên}
\label{algo:fraud}
\KwIn{Tập bản ghi điểm danh $A$ của lớp}
\KwOut{Danh sách mẫu hành vi đáng ngờ $S$}
$S \leftarrow \emptyset$\;
\ForEach{sinh viên $s$ trong $A$}{
  \lIf{tỷ lệ xác minh cùng một nhóm bạn $\geq 80\%$}{thêm \textit{always\_same\_peers}}
  \lIf{khoảng cách giữa hai lượt xác minh $< 30$ giây}{thêm \textit{rapid\_verification}}
  \lIf{số lượt ngang hàng $< 3$ ở quá nửa số phiên}{thêm \textit{low\_peer\_count}}
  \lIf{khuôn mặt không khớp hoặc độ tin cậy $< 0{,}3$}{thêm \textit{face\_mismatch}}
}
\Return $S$\;
\end{algorithm}

\textbf{Kiểm tra vị trí GPS (hàng rào địa lý).} Khi giảng viên mở phiên điểm danh, ứng dụng
ghi lại toạ độ GPS của thiết bị giảng viên làm tâm khu vực lớp học cùng bán kính cho phép
(mặc định 200~m). Khi sinh viên điểm danh, thiết bị lấy toạ độ hiện tại và tính khoảng cách
tới tâm bằng công thức Haversine; nếu vượt quá bán kính, lượt điểm danh bị từ chối. Trường
hợp sinh viên tắt định vị hoặc từ chối quyền, hệ thống không chặn cứng mà đánh dấu bản ghi
``cần xem xét'' để giảng viên duyệt. Việc kiểm tra được thực hiện đồng thời ở thiết bị và
trên máy chủ (Cloud Function tính lại khoảng cách), nên không thể vượt qua bằng cách can
thiệp phía máy khách.

\subsection{Tính năng Cổng máy chiếu (Teacher Display)}
\label{sec:projector}

Mã QR hiển thị trên màn hình điện thoại của giảng viên thường quá nhỏ đối với các lớp học
đông. Để khắc phục, hệ thống thiết kế tính năng \textit{Cổng máy chiếu}: trang quản trị
web mở trên máy chiếu sẽ sinh một mã ghép nối (pairing token) ngẫu nhiên dài 128 bit, có
thời hạn 60 giây. Giảng viên dùng Mini App quét mã ghép nối này để liên kết phiên điểm
danh với màn hình máy chiếu; sau đó mã QR động cỡ lớn sẽ được hiển thị trên máy chiếu.

Trạng thái của mã ghép nối chỉ chuyển một chiều từ \texttt{pending} sang \texttt{paired}
và không thể quay lại hay bị ghép nối hai lần, nhằm chống tấn công chiếm quyền hiển thị.
Mã QR trên máy chiếu và trên Mini App được đồng bộ theo cùng một đồng hồ để bảo đảm tính
nhất quán khi làm mới. Trình tự ghép nối được minh hoạ trong Hình~\ref{fig:seq-pairing}.

\begin{figure}[H]
  \centering
  \includegraphics[width=0.8\textwidth]{seq_pairing}
  \caption{Lược đồ tuần tự ghép nối Cổng máy chiếu}
  \label{fig:seq-pairing}
\end{figure}

Hình~\ref{fig:projector-ui} thể hiện giao diện trang Cổng máy chiếu khi chạy trên
trình duyệt.

\begin{figure}[H]
  \centering
  \includegraphics[width=0.95\textwidth]{admin_projector}
  \caption{Giao diện Cổng máy chiếu trên trình duyệt}
  \label{fig:projector-ui}
\end{figure}

% ═══════════════════════════════════════════════════════════════
% PHẦN B — TRIỂN KHAI  (Đợt 2)
% ═══════════════════════════════════════════════════════════════

\section{Xây dựng ứng dụng}

Phần này trình bày quá trình hiện thực hoá hệ thống thành ba thành phần: Mini App, trang
quản trị và các hàm xử lý phía máy chủ.

\subsection{Triển khai Mini App}
\label{sec:trien-khai-miniapp}

Phần Mini App được hiện thực bằng React và TypeScript với quy mô khoảng 15.900 dòng mã,
tổ chức thành 21 trang giao diện, 27 thành phần dùng chung, 10 hook tùy biến và 16 lớp
dịch vụ (service) đảm nhiệm việc giao tiếp với Firestore và Cloud Functions.

Mối quan hệ phân lớp giữa các thành phần giao diện, hook, dịch vụ và Firebase được minh hoạ
trong Hình~\ref{fig:class-layer}.

\begin{figure}[H]
  \centering
  \includegraphics[width=0.9\textwidth]{class_layer}
  \caption{Sơ đồ phân lớp các thành phần phần mềm của Mini App}
  \label{fig:class-layer}
\end{figure}

Các trang được phân theo vai trò người dùng:
\begin{itemize}[leftmargin=1.5em]
  \item \textbf{Chung:} màn hình khởi động, giới thiệu, đăng nhập, trang chủ, hồ sơ và
  trợ lý hỏi đáp.
  \item \textbf{Sinh viên:} danh sách lớp, lịch học, đăng ký khuôn mặt, màn hình điểm
  danh bốn bước và lịch sử điểm danh.
  \item \textbf{Giảng viên:} danh sách lớp, chi tiết lớp, tạo và quản lý phiên điểm danh,
  giám sát theo thời gian thực, duyệt kết quả, báo cáo gian lận và thống kê.
\end{itemize}

Việc giao tiếp với nền tảng Zalo được thực hiện qua ZMP SDK, gồm xác thực người dùng,
truy cập máy ảnh để quét mã QR và chụp ảnh khuôn mặt, và truy cập vị trí GPS. Một số màn
hình tiêu biểu được minh hoạ trong Hình~\ref{fig:miniapp}.

\begin{figure}[H]
  \centering
  \begin{subfigure}[t]{0.30\textwidth}\centering\includegraphics[width=\linewidth]{mini_home}\caption{Trang chủ}\end{subfigure}\hfill
  \begin{subfigure}[t]{0.30\textwidth}\centering\includegraphics[width=\linewidth]{mini_classes}\caption{Danh sách lớp}\end{subfigure}\hfill
  \begin{subfigure}[t]{0.30\textwidth}\centering\includegraphics[width=\linewidth]{mini_schedule}\caption{Thời khoá biểu}\end{subfigure}

  \vspace{0.4cm}
  \begin{subfigure}[t]{0.30\textwidth}\centering\includegraphics[width=\linewidth]{mini_search}\caption{Tra cứu}\end{subfigure}\hfill
  \begin{subfigure}[t]{0.30\textwidth}\centering\includegraphics[width=\linewidth]{mini_absence}\caption{Đơn xin nghỉ}\end{subfigure}\hfill
  \begin{subfigure}[t]{0.30\textwidth}\centering\includegraphics[width=\linewidth]{mini_profile}\caption{Hồ sơ}\end{subfigure}
  \caption{Một số màn hình chính của Mini App}
  \label{fig:miniapp}
\end{figure}

\subsection{Triển khai trang quản trị}
\label{sec:trien-khai-admin}

Trang quản trị là một ứng dụng web độc lập dành cho phòng đào tạo, được xây dựng bằng
React kết hợp thư viện giao diện Ant Design, thư viện biểu đồ Recharts và thư viện
\texttt{xlsx} để nhập/xuất dữ liệu Excel. Trang gồm 12 màn hình chính, trong đó nổi bật
là: bảng điều khiển thống kê tổng quan, quản lý người dùng, quản lý lớp học, theo dõi
điểm danh, duyệt đơn xin nghỉ, báo cáo gian lận và quản lý mã mời giảng viên.

Hai tính năng đặc thù của trang quản trị là \textbf{nhập danh sách sinh viên và lớp hàng
loạt từ tệp Excel} và \textbf{xuất báo cáo điểm danh ra Excel}, giúp giảm đáng kể công
sức nhập liệu thủ công cho phòng đào tạo. Trang quản trị được triển khai trên Firebase
Hosting. Giao diện được minh hoạ trong Hình~\ref{fig:admin}.

\begin{figure}[H]
  \centering
  \includegraphics[width=0.95\textwidth]{admin_dashboard}
  \caption{Bảng điều khiển tổng quan của trang quản trị}
  \label{fig:admin}
\end{figure}

Quy trình xử lý đơn xin nghỉ giữa sinh viên và quản trị viên được minh hoạ trong
Hình~\ref{fig:seq-absence}.

\begin{figure}[H]
  \centering
  \includegraphics[width=0.75\textwidth]{seq_absence}
  \caption{Lược đồ tuần tự xử lý đơn xin nghỉ}
  \label{fig:seq-absence}
\end{figure}

\subsection{Triển khai Cloud Functions}
\label{sec:trien-khai-functions}

Toàn bộ nghiệp vụ nhạy cảm được xử lý phía máy chủ thông qua 24 hàm Cloud Functions, được
phân nhóm theo chức năng như trong Bảng~\ref{tab:functions}. Việc tập trung logic ở máy
chủ bảo đảm khoá bí mật HMAC, thuật toán chống gian lận và việc gán vai trò không bị can
thiệp từ phía thiết bị người dùng.

\begin{table}[H]
\centering
\small
\caption{Phân nhóm các hàm Cloud Functions}
\label{tab:functions}
\begin{tabular}{|p{3.2cm}|p{10cm}|}
\hline
\textbf{Nhóm} & \textbf{Các hàm} \\
\hline
Phiên điểm danh & \texttt{startSession}, \texttt{endSession} \\ \hline
Điểm danh & \texttt{scanTeacher}, \texttt{scanPeer}, \texttt{submitFaceResult}, \texttt{manualAttendance}, \texttt{reviewAttendance}, \texttt{calculateTrustPolicies} \\ \hline
Khuôn mặt & \texttt{registerFace}, \texttt{verifyFace} \\ \hline
Chống gian lận & \texttt{analyzeFraud}, \texttt{weeklyFraudAnalysis} \\ \hline
Xác thực \& phân quyền & \texttt{createFirebaseToken}, \texttt{verifyZaloToken}, \texttt{resolveZaloPhoneToken}, \texttt{assignTeacherRole} \\ \hline
Quản trị & \texttt{adminDashboardStats}, \texttt{adminBulkAssignStudents}, \texttt{adminReviewAbsenceRequest}, \texttt{adminCreateAccount} \\ \hline
Thông báo \& AI & \texttt{notifySessionStarted}, \texttt{aiChat} \\ \hline
\end{tabular}
\end{table}

Một số hàm tiêu biểu được mô tả chi tiết hơn như sau. Nhóm \textbf{phiên điểm danh} gồm
\texttt{startSession} (sinh khoá bí mật HMAC và khởi tạo phiên) và \texttt{endSession} (đóng
phiên, chốt kết quả). Nhóm \textbf{điểm danh} là nhóm cốt lõi: \texttt{scanTeacher} xác minh
mã QR của giảng viên, \texttt{scanPeer} xác minh mã ngang hàng giữa các sinh viên,
\texttt{submitFaceResult} ghi nhận kết quả khuôn mặt, còn \texttt{calculateTrustPolicies} tổng
hợp và tính điểm tin cậy cho toàn phiên.

Nhóm \textbf{xác thực và phân quyền} đảm nhiệm việc bắc cầu danh tính: \texttt{verifyZaloToken}
kiểm tra token Zalo, \texttt{createFirebaseToken} sinh Custom Token, \texttt{assignTeacherRole}
gán vai trò giảng viên (chỉ ở phía máy chủ). Nhóm
\textbf{chống gian lận} gồm \texttt{analyzeFraud} chạy theo yêu cầu và
\texttt{weeklyFraudAnalysis} chạy định kỳ. Việc tập trung toàn bộ các nghiệp vụ nhạy cảm tại
tầng máy chủ là một quyết định thiết kế quan trọng nhằm bảo đảm an toàn cho hệ thống. Mã nguồn
chi tiết của một số hàm cốt lõi được trích dẫn trong Phụ lục~\ref{chapter:appendixC}.

% ═══════════════════════════════════════════════════════════════
% PHẦN C — KIỂM THỬ & ĐÁNH GIÁ  (Đợt 3)
% ═══════════════════════════════════════════════════════════════

\section{Kiểm thử và đánh giá}

Phần này trình bày kết quả kiểm thử cùng các đánh giá về bảo mật và hiệu năng của hệ thống.

\subsection{Kiểm thử}
\label{sec:kiem-thu}

Do hệ thống bao gồm nhiều thành phần tương tác trực tiếp với phần cứng thiết bị (máy ảnh,
cảm biến vị trí) và với nền tảng Zalo, việc kiểm thử được thực hiện chủ yếu theo hình thức
\textbf{kiểm thử thủ công} theo hướng hộp đen (black-box) trên thiết bị thật. Mục tiêu là
kiểm chứng tính đúng đắn của từng luồng nghiệp vụ từ góc nhìn người dùng cuối, thay vì chỉ
kiểm tra từng hàm riêng lẻ. Bên cạnh đó, toàn bộ mã nguồn được kiểm tra kiểu tĩnh bằng trình
biên dịch TypeScript (\texttt{tsc -{}-noEmit}) nhằm phát hiện sớm các lỗi kiểu dữ liệu trước
khi chạy.

Môi trường kiểm thử sử dụng bốn điện thoại thông minh cài sẵn Zalo: một máy đóng vai giảng
viên và ba máy đóng vai sinh viên (để mô phỏng bước xác minh ngang hàng vốn cần tối thiểu ba
lượt), cùng kết nối Internet qua Wi-Fi hoặc mạng di động. Ngoài ra, ứng dụng còn cung cấp một
chế độ giả lập dữ liệu (mock mode) cho phép kiểm thử riêng lẻ từng bước trong luồng điểm danh
mà không phụ thuộc vào dữ liệu thật.

Bộ kịch bản kiểm thử gồm \textbf{72 kịch bản} (TC-01 đến TC-72) được tổ chức thành 19 nhóm
chức năng, bao phủ toàn bộ các luồng chính của ứng dụng và 17 tình huống biên. Mỗi kịch bản
mô tả rõ điều kiện tiền đề, các bước thực hiện, kết quả mong đợi và ô đánh dấu Đạt/Không đạt.
Phân bố số lượng kịch bản theo nhóm chức năng được trình bày trong
Bảng~\ref{tab:test-coverage}.

\begin{table}[H]
\centering
\small
\caption{Phân bố kịch bản kiểm thử theo nhóm chức năng}
\label{tab:test-coverage}
\begin{tabular}{|p{9cm}|c|}
\hline
\textbf{Nhóm chức năng} & \textbf{Số kịch bản} \\
\hline
Xác thực \& onboarding & 6 \\ \hline
Trang chủ (dashboard) & 3 \\ \hline
Quản lý lớp học (giảng viên) & 3 \\ \hline
Tham gia lớp học (sinh viên) & 4 \\ \hline
Phiên điểm danh (giảng viên) & 5 \\ \hline
Điểm danh bốn bước (sinh viên) & 4 \\ \hline
Theo dõi thời gian thực & 3 \\ \hline
Kết thúc phiên \& điểm tin cậy & 2 \\ \hline
Xem xét \& ghi đè kết quả & 4 \\ \hline
Điểm danh thủ công & 2 \\ \hline
Lịch sử điểm danh (sinh viên) & 2 \\ \hline
Đăng ký khuôn mặt & 2 \\ \hline
Hồ sơ người dùng & 4 \\ \hline
Trợ lý AI & 4 \\ \hline
Báo cáo gian lận & 2 \\ \hline
Thống kê (analytics) & 1 \\ \hline
Cấu hình bước tuỳ chọn & 3 \\ \hline
Lịch học & 1 \\ \hline
Tình huống biên (edge case) & 17 \\ \hline
\textbf{Tổng cộng} & \textbf{72} \\ \hline
\end{tabular}
\end{table}

Một số kịch bản tiêu biểu, đại diện cho các cơ chế cốt lõi của hệ thống, được liệt kê trong
Bảng~\ref{tab:test-sample}.

\begin{table}[H]
\centering
\small
\caption{Một số kịch bản kiểm thử tiêu biểu}
\label{tab:test-sample}
\begin{tabular}{|l|p{5.8cm}|p{5.8cm}|}
\hline
\textbf{Mã} & \textbf{Mục tiêu kiểm thử} & \textbf{Kết quả mong đợi} \\
\hline
TC-05 & Ràng buộc định dạng MSSV & MSSV không đủ 8 số hoặc không bắt đầu bằng ``20'' bị từ chối \\ \hline
TC-18 & Mã QR hiển thị và xoay 30 giây & Mã QR tự làm mới sau mỗi 30 giây \\ \hline
TC-23 & Xác minh khuôn mặt tự động & So khớp đạt khi độ tin cậy không nhỏ hơn 0{,}7 \\ \hline
TC-24 & Trao đổi mã QR ngang hàng & Yêu cầu tối thiểu ba lượt xác minh hợp lệ \\ \hline
TC-30 & Tính điểm tin cậy & Trả về present/review/absent đúng quy tắc \\ \hline
TC-56 & Mã QR hết hạn & Mã quá 90 giây bị từ chối \\ \hline
TC-57 & Tự quét mã của chính mình & Hệ thống chặn việc tự xác minh ngang hàng \\ \hline
\end{tabular}
\end{table}

Nhóm tình huống biên (17 kịch bản) tập trung vào các trường hợp dễ phát sinh lỗi hoặc bị lợi
dụng, ví dụ: mã QR đã hết hạn, sinh viên tự quét mã của chính mình ở bước ngang hàng, quét
trùng một bạn đã xác minh, hoặc mất kết nối mạng giữa chừng. Đây cũng chính là những điểm mà
cơ chế chống gian lận và khả năng chống chịu lỗi của hệ thống được kiểm chứng.

Quá trình kiểm thử được thực hiện trên thiết bị thật, mỗi kịch bản được đánh dấu Đạt hoặc
Không đạt trong tài liệu kế hoạch kiểm thử. Các luồng nghiệp vụ chính --- đăng nhập và phân
vai, tạo và quản lý phiên điểm danh, điểm danh bốn bước, tính điểm tin cậy, xem xét và ghi đè
kết quả --- đều hoạt động đúng kỳ vọng. Một số lỗi phát hiện trong quá trình kiểm thử, chủ yếu
liên quan đến việc xử lý tình huống biên và độ trễ khởi động máy ảnh, đã được ghi nhận và khắc
phục.

\subsection{Đánh giá bảo mật}
\label{sec:bao-mat}

Vì điểm danh chống gian lận về bản chất là một bài toán bảo mật, hệ thống được rà soát một
cách có hệ thống theo các nhóm chủ đề quen thuộc (tương tự phân loại của OWASP): xác thực,
phân quyền, kiểm tra đầu vào, mã hoá, bảo vệ dữ liệu, an toàn của tầng máy chủ và xử lý lỗi.
Mỗi lĩnh vực được đánh giá cả về những gì đã đạt được lẫn những gì cần cải thiện, kèm theo một
mức điểm định tính. Kết quả tổng hợp được trình bày trong Bảng~\ref{tab:security-score}.

\begin{table}[H]
\centering
\small
\caption{Bảng điểm đánh giá bảo mật theo lĩnh vực}
\label{tab:security-score}
\begin{tabular}{|p{5.2cm}|p{6.4cm}|c|}
\hline
\textbf{Lĩnh vực} & \textbf{Hiện trạng} & \textbf{Điểm} \\
\hline
Xác thực & Đã có cơ chế cơ bản qua Zalo và Custom Token & 7/10 \\ \hline
Phân quyền & Firestore rules chặt chẽ, giới hạn vai trò & 9/10 \\ \hline
Kiểm tra đầu vào & Nhiều hàm kiểm tra và làm sạch dữ liệu & 8/10 \\ \hline
Mã hoá mã QR & HMAC-SHA256 mạnh, khoá bí mật được bảo vệ & 9/10 \\ \hline
Chống gian lận & Kiểm tra phía máy chủ & 7/10 \\ \hline
Bảo vệ khoá API & Đã chuyển toàn bộ về phía máy chủ & 8/10 \\ \hline
Bảo vệ dữ liệu & Rules tốt, song localStorage chưa mã hoá & 7/10 \\ \hline
Cloud Functions & Có xác thực và giới hạn tần suất & 8/10 \\ \hline
Xử lý lỗi & Có cơ chế bắt lỗi và thông báo & 7/10 \\ \hline
Mạng \& ngoại tuyến & Có hàng đợi ngoại tuyến & 6/10 \\ \hline
\multicolumn{2}{|r|}{\textbf{Tổng điểm}} & \textbf{7{,}6/10} \\ \hline
\end{tabular}
\end{table}

Trong quá trình rà soát, sáu lỗ hổng đã được phát hiện và khắc phục, nâng điểm bảo mật tổng
thể từ 6{,}2/10 lên 7{,}6/10. Các lỗ hổng cùng biện pháp khắc phục tương ứng được tóm tắt
trong Bảng~\ref{tab:security-fix}.

\begin{table}[H]
\centering
\small
\caption{Các lỗ hổng bảo mật đã phát hiện và khắc phục}
\label{tab:security-fix}
\begin{tabular}{|p{4.6cm}|c|p{6.6cm}|}
\hline
\textbf{Lỗ hổng} & \textbf{Mức độ} & \textbf{Biện pháp khắc phục} \\
\hline
Khoá API lộ trong gói mã phía trình duyệt & Cao & Chuyển lời gọi sang Cloud Function, xoá khoá khỏi phần frontend \\ \hline
Khoá API thật nằm trong tệp ví dụ cấu hình & Cao & Thay bằng giá trị mẫu và thu hồi khoá cũ \\ \hline
Khoá bí mật HMAC lộ qua Firestore & Cao & Chuyển vào subcollection chỉ giảng viên đọc; Cloud Function kiểm tra mã \\ \hline
Không giới hạn vai trò khi tạo người dùng & Cao & Rule chỉ cho phép vai trò rỗng hoặc ``student'' \\ \hline
Đường dự phòng điểm danh bỏ qua kiểm tra & Trung bình & Siết chặt rule và bắt buộc đi qua Cloud Function \\ \hline
Giả mạo xác minh ngang hàng phía thiết bị & Trung bình & Chuyển khâu xác minh sang máy chủ (hàm \texttt{scanPeer}) \\ \hline
\end{tabular}
\end{table}

Nhìn chung, hệ thống có nền tảng bảo mật tốt: mã QR được ký bằng HMAC-SHA256 với khoá bí mật
được cô lập, quy tắc Firestore ngăn người dùng tự nâng quyền, toàn bộ nghiệp vụ nhạy cảm được
xử lý phía máy chủ, và mã nguồn không sử dụng các hàm nguy hiểm như \texttt{eval} hay
\texttt{innerHTML}. Một số hạn chế vẫn còn ở mức ưu tiên trung bình và dài hạn: dữ liệu trong
\texttt{localStorage} chưa được mã hoá; cơ chế giới hạn tần suất và chống phát lại (nonce) hiện
lưu trong bộ nhớ tạm nên không bền vững khi mở rộng quy mô; và hệ thống chưa có nhật ký kiểm
toán (audit log) hay tự đăng xuất sau thời gian dài. Đây là các định hướng để nâng mức độ bảo
mật lên cao hơn.

\subsection{Đánh giá hiệu năng}
\label{sec:hieu-nang}

Hiệu năng phía thiết bị được tối ưu bằng cách \textbf{tải trễ} (lazy-load) các mô hình nhận
diện khuôn mặt: bộ mô hình có dung lượng khoảng 6{,}4~MB chỉ được tải khi người dùng vào các
trang cần đến, thay vì tải ngay lúc khởi động, nhờ đó thời gian mở ứng dụng được giữ ở mức
thấp. Hệ thống cũng lựa chọn mô hình phát hiện khuôn mặt gọn nhẹ để có thể chạy trực tiếp trên
thiết bị di động.

Tính năng giám sát điểm danh theo thời gian thực dựa trên cơ chế lắng nghe thay đổi
(\texttt{onSnapshot}) của Firestore, cho phép cập nhật gần như tức thời mà không cần liên tục
hỏi vòng máy chủ (polling), qua đó tiết kiệm băng thông và phản hồi nhanh. Ở phía máy chủ, các
hàm Cloud Functions hoạt động theo mô hình serverless tự co giãn theo tải và được đặt tại khu
vực \texttt{asia-southeast1} (gần người dùng tại Việt Nam) nhằm giảm độ trễ. Khi mạng chập
chờn, hàng đợi ngoại tuyến và cơ chế lưu đệm cục bộ giúp người dùng vẫn thao tác được, dữ liệu
sẽ được đồng bộ lại khi có kết nối.

Về hạn chế, độ trễ khởi động nguội (cold start) của Cloud Functions có thể làm chậm các yêu
cầu đầu tiên sau một khoảng thời gian không hoạt động; bên cạnh đó, hiệu năng của khâu nhận
diện khuôn mặt phụ thuộc vào cấu hình của từng thiết bị. Việc đo đạc hiệu năng một cách định
lượng (thời gian phản hồi, thông lượng) trên nhiều cấu hình thiết bị khác nhau là một hướng
phát triển tiếp theo của đồ án.

\section{Triển khai và phát hành}
\label{sec:phat-hanh}

Hệ thống được phát hành theo nhiều kênh tương ứng với từng thành phần. Mini App được đóng gói
và phát hành lên nền tảng Zalo thông qua công cụ dòng lệnh (\texttt{zmp deploy}) và phải trải
qua quy trình kiểm duyệt của Zalo trước khi đến tay người dùng. Trang quản trị được triển khai
trên Firebase Hosting; còn các hàm Cloud Functions, quy tắc bảo mật và chỉ mục của Firestore
được triển khai qua công cụ Firebase CLI. Ở giai đoạn hiện tại, hệ thống chạy trên gói miễn
phí (Spark); để đưa vào sử dụng thực tế ở quy mô lớn, cần nâng lên gói trả phí (Blaze) nhằm
khai thác đầy đủ năng lực của Cloud Functions. Tô-pô triển khai vật lý của hệ thống được minh
hoạ trong Hình~\ref{fig:deployment}.

\begin{figure}[H]
  \centering
  \includegraphics[width=0.95\textwidth]{deployment}
  \caption{Sơ đồ triển khai vật lý của hệ thống}
  \label{fig:deployment}
\end{figure}

Về tuân thủ chính sách nền tảng, ứng dụng tuân theo quy định của Zalo Mini App: không yêu cầu
cấp quyền ngay khi mở ứng dụng mà chỉ xin quyền khi người dùng chủ động thực hiện thao tác
liên quan; định danh người dùng được lấy qua \texttt{getUserID} (không cần quyền), trong khi
thông tin cá nhân và số điện thoại được xin theo từng phạm vi riêng. Ứng dụng cũng cung cấp
đầy đủ Điều khoản sử dụng và Chính sách quyền riêng tư cho người dùng.

Về bảo vệ dữ liệu cá nhân, dữ liệu khuôn mặt chỉ được lưu dưới dạng véc-tơ đặc trưng 128 chiều
chứ không lưu ảnh gốc trên máy chủ, giúp giảm rủi ro lộ lọt dữ liệu sinh trắc học. Trong quá
trình chuẩn bị phát hành rộng rãi, sản phẩm cũng được điều chỉnh theo hướng trung lập về
thương hiệu --- gỡ bỏ các tham chiếu định danh của một cơ sở đào tạo cụ thể khỏi giao diện và
chính sách --- qua đó thuận lợi hơn cho việc triển khai cho nhiều đơn vị khác nhau.

\subsection*{Kết chương}

Chương này đã trình bày trọn vẹn quá trình hiện thực hoá hệ thống: từ thiết kế kiến trúc ba
tầng, cơ sở dữ liệu Firestore, luồng xác thực --- phân quyền và các thuật toán cốt lõi (mã QR
HMAC xoay, điểm danh bốn bước, chống gian lận nhiều lớp, cổng máy chiếu), đến việc triển khai
thành ba sản phẩm thực tế (Mini App, trang quản trị và các hàm máy chủ), và cuối cùng là kết
quả kiểm thử cùng các đánh giá về bảo mật (đạt 7{,}6/10), hiệu năng và quy trình phát hành.
Kết quả cho thấy hệ thống đã đáp ứng được các yêu cầu chức năng và phi chức năng đã đề ra,
đồng thời cũng bộc lộ một số hạn chế làm tiền đề cho các cải tiến tiếp theo. Chương sau sẽ đi
sâu vào những giải pháp và đóng góp nổi bật nhất của đồ án.

\end{document}
